\section{Introduction}
\label{sec:intro}

Substandard and falsified medicines harm patients in every region, and low- and middle-income countries carry most of the burden \cite{who2017oneinten,ozawa2018prevalence}. In India, drug samples are tested by central and state laboratories, and CDSCO publishes the batches that fail as monthly alerts \cite{cdsco2026portal}. The volume of these alerts has grown quickly: the portal lists \RowsYearFirst{} entries for 2019 and \RowsYearTwentyFive{} for 2025, with \RowsYearTwentySix{} more in the first nine months of 2026.

The alerts are published as tables of product names, batch numbers, dates and manufacturers. Someone who wants to know whether the strip in their hand is affected has to find the right month, read the batch number from the foil and compare it by hand. Packaging codes do not help here. Since August 2023 the top 300 brands carry a QR code under Schedule H2 \cite{cdsco2023faqqr}, and such codes, like serialization schemes elsewhere \cite{who2021traceability,fda2026dscsaexemptions}, show where a pack came from. WHO notes that verifying a unique identifier cannot check the medicine inside the package \cite{who2021traceability}. A batch that a licensed manufacturer made badly still carries valid codes. Only the laboratory result says that the batch failed, and the alert list is where that result is published.

Automating the check looks like a lookup, but a lookup script fails in two ways. The first is reading. OCR on blister foil makes confident single-character errors: in our tests a scuffed \texttt{B} was read as \texttt{3} with a reported confidence of 1.0, and an enhanced re-read returned the same answer. A misread batch number either misses a listed batch or, worse, lands on another one. The second is matching. Batch numbers are short and are chosen by each manufacturer, so the same string appears on unrelated products. In the CDSCO data, \CollisionKeys{} batch keys are shared by unrelated products, and \SiblingKeys{} keys have a listed neighbour from the same manufacturer that differs in one easily confused character, for example \texttt{R21119} and \texttt{R21719}. A confident misread can therefore land exactly on a listed sibling, and every other field on the pack (product, maker, dates) will agree with it.

StripCheck treats the check as a decision under uncertainty with a person in the loop. A bounded agent reads the pack (QR code, OCR, an enhanced re-read and, optionally, a small vision-language model as an independent second reader), retrieves alert entries by exact and confusion-aware batch search, and grades each candidate field by field, which is the structure of blocked record linkage \cite{fellegi1969theory}. When one character is in doubt it shows the person the camera's crop of that character and asks. Before it raises any alarm it requires the batch to come from a QR code, from the person, or from two independent reads that agree. A verifier agent checks every draft verdict against hard rules and downgrades it to ``cannot verify'' when the evidence does not support it. A small local language model plans at branch points and judges ambiguous product names, with a deterministic rule planner as fallback, so the system runs on a laptop with no paid services and no network access during a check.

Open tools built on CDSCO's alerts appeared in 2025 and 2026: one looks up typed batch numbers, one reads pack photos and declines to answer when it sees look-alike characters, and one matches alerts against distributor stock (Section~\ref{sec:related}). StripCheck differs in asking the person to settle the doubtful character, in corroborating every field of a photographed consumer pack before any alarm, and in running entirely offline.

The contributions of this work are:
\begin{itemize}
\item an ingest agent for CDSCO's alert portal and legacy PDF lists, with validation, grounded repair and quarantine, and an analysis of \NumMonths{} months of alert data that measures batch collisions, confusable sibling batches and how long flagged batches remain in circulation;
\item an agent design that combines confusion-aware retrieval, record-linkage corroboration, a confirmation gate before alarms and a rule-based verifier with abstention, and that addresses the sibling-misread false alarm, which a plain lookup cannot detect;
\item an evaluation protocol with synthetic packs carrying calibrated misreads and sibling traps, a degradation ladder, printed and photographed packs, real strips, ablations and preregistered decision rules against a lookup baseline;
\item an offline implementation with a phone-friendly interface, released for reproduction \todo{repository link}.
\end{itemize}
